\subsection{Servidor}


En esta sección se presenta la selección del broker MQTT y del servidor para el Módulo 4
del proyecto, que integra Python con Flask, MongoDB, React y los microcontroladores para gestionar
datos de fermentación en tiempo real. La elección se fundamenta en los requerimientos de escalabilidad
horizontal, flexibilidad para datos dinámicos JSON, baja latencia en la comunicación MQTT y 
optimización de costos para un despliegue inicial en un sistema IoT compacto. 
A ontinuación se evalúan opciones 
de brokers MQTT y servicios en la nube, considerando su compatibilidad con la arquitectura del 
sistema, facilidad de integración con Flask y React, y capacidad para soportar análisis y
visualización de datos en tiempo real. Por consiguiente, se seleccionan las soluciones que
mejor equilibran rendimiento, simplicidad y viabilidad económica.

\subsubsection{Servidor en la nube o local}
De la tabla  \cref{tab:serlocal} se selecciona un servidor de forma local

como la mejor opción, y apesar de que un servidor en la nube ofrece ventajas significativas,
las limitaiciones de costos y la dependencia de la conectividad a internet hacen que un servidor local 
sea más adecuado para este proyecto,sin embargo las herramientas previamente seleccionadas permiten 
poder escalar hacia una arquitectura en la nube en un futuro.Aun así, se sigue cumpliendo los
requerimientos funcionales \ref{it:RF12} y no funcionales RNF04, al mantener 
la capacidad de almacenar y procesar datos de manera eficiente, garantizando la disponibilidad .

\begingroup
\setlength{\tabcolsep}{5pt} 
\renewcommand{\arraystretch}{1.5} 
\small 

% Usamos longtable para evitar saltos de página innecesarios
\begin{longtable}{|p{2.6cm}|>{\RaggedRight}p{5.8cm}|>{\RaggedRight}p{5.8cm}|}
    \caption{Comparación de Servidor en la nube o local.} \label{tab:serlocal} \\
    \hline
    \textbf{Categoría} & \textbf{Servidor Local} \cite{zotero-item-250} & \textbf{Servidor en la Nube} \cite{2017} \\ \hline
    \endfirsthead
    
    \multicolumn{3}{l}{\footnotesize\itshape -- continuación de la página anterior} \\ \hline
    \textbf{Categoría} & \textbf{Servidor Local} & \textbf{Servidor en la Nube} \\ \hline
    \endhead

    Costo & Sin costos recurrentes, solo inversión inicial en hardware. & Costos mensuales bajo demanda (\$10-40/mes para proyectos pequeños). \\ \hline
    
    Escalabilidad & Limitada, requiere hardware adicional y configuración manual. & Alta, con autoescalado y balanceo de carga automáticos. \\ \hline
    
    Mantenimiento & Responsabilidad del equipo, incluye actualizaciones y backups. & Gestionado por el proveedor, minimizando carga operativa. \\ \hline
    
    Disponibilidad & Vulnerable a fallos de hardware o energía local. & Alta disponibilidad (99.9\%+ SLA) con redundancia. \\ \hline
    
    Integración & Configuración manual para MQTT, Flask, MongoDB, React. & Integración nativa con servicios como MongoDB Atlas y APIs cloud. \\ \hline
    
    Latencia & Muy baja en redes locales, ideal para pruebas iniciales. & Depende de internet, optimizable mediante regiones cercanas. \\ \hline
\end{longtable}
\endgroup


\subsubsection{MQTT Servidor}
Basado en la tabla comparativa \cref{tab:serMQTT}, se selecciona \textbf{EMQX Cloud}
para el Módulo 1, dado que su escalabilidad horizontal, soportada por clustering distribuido,
permite conexiones simultáneas, alineándose con el requerimiento RNF01 
de añadir dispositivos sin afectar los existentes. Su soporte para MQTT 5.0 y análisis en 
tiempo real optimiza el procesamiento de datos de fermentación enviados por los microcontroladores, mientras
que su integración nativa con Python mediante \texttt{paho-mqtt} y JSON asegura un manejo 
eficiente en Flask y MongoDB. Además, su alta disponibilidad , 
seguridad avanzada con TLS/ACLS y despliegue en regiones de AWS, GCP o Azure garantizan
fiabilidad y baja latencia global para visualizaciones en React. Aunque implica costos,
su modelo ``pay as you go'' es 
asequible para proyectos pequeños, y su documentación extensa facilita la transición
desde entornos locales. En contraste, HiveMQ Cloud presenta costos más altos
(\textasciitilde\$1,000/mes), NanoMQ y VerneMQ requieren mayor configuración técnica,
y AWS IoT Core introduce dependencia de su ecosistema. 
Aunque Eclipse Mosquitto es ligero y gratuito, su escalabilidad limitada 
y dependencia de hardware lo hacen menos viable para la nube. Por consiguiente,
EMQX Cloud equilibra rendimiento, simplicidad y viabilidad económica para este sistema
IoT en la nube.



\begingroup
\setlength{\tabcolsep}{4pt} % Espacio entre columnas
\renewcommand{\arraystretch}{1.5} % Altura de las filas
\small % Fuente pequeña para aprovechar el ancho

% Se usa longtable en lugar de tabularx. Se define p{ancho} para columnas fijas 
% o >{\RaggedRight}X para columnas flexibles si usas tabularx.
% Aquí usaremos p{ancho} para una mayor estabilidad y \textwidth en longtable.
\begin{longtable}{|p{2.1cm}|p{3.5cm}|p{3cm}|p{2.8cm}|p{2.8cm}|}
\caption{Comparación de servicios MQTT.}
\label{tab:serMQTT}
\\
\hline
\textbf{Broker MQTT} & \textbf{Descripción} & \textbf{Ventaja} & \textbf{Desventaja} & \textbf{Costos} \\
\hline
\endfirsthead
% --- Encabezado que se repite en la segunda página ---
\multicolumn{5}{l}{\small\bfseries \tablename\ \thetable} \\
\hline
\textbf{Broker MQTT} & \textbf{Descripción} & \textbf{Ventaja} & \textbf{Desventaja} & \textbf{Costos} \\
\hline
\endhead
\hline
\multicolumn{5}{r}{\small\itshape } \\
\endfoot
\hline
\endlastfoot

\textbf{Eclipse Mosquitto}  \cite{zotero-item-236}  & Broker ligero de código abierto, soporta QoS 0-2, TLS y ACLs, ideal para despliegues locales. & Ligero y fácil de configurar, ideal para proyectos pequeños con baja latencia. & Escalabilidad limitada ya que depende mucho de los recursos del sistema. & Gratuito (licencia EPL/EDL). \\
\hline
\textbf{EMQX} \cite{zotero-item-228} & Broker escalable con soporte MQTT 5.0, clustering distribuido y análisis en tiempo real. & Alta escalabilidad horizontal, soporta millones de conexiones con clustering. & Configuración compleja para entornos pequeños. & Open Source: Gratuito; Enterprise: Desde \textasciitilde\$0.10/conexión/mes (\textasciitilde\$100/mes plan básico). \\
\hline
\textbf{VerneMQ} \cite{2025d}  & Broker basado en Erlang, enfocado en escalabilidad masiva y autenticación avanzada. & Escalabilidad masiva con clustering basado en Erlang, ideal para IoT. & Requiere experiencia técnica para optimización y mantenimiento. & Gratuito; soporte enterprise \textasciitilde\$5,000/año. \\
\hline
\textbf{HiveMQ} \cite{zotero-item-238} & Broker empresarial con soporte MQTT 5.0, seguridad RBAC y compatibilidad con Kafka. & Seguridad avanzada y soporte MQTT 5.0, optimizado para entornos enterprise. & Costo elevado para despliegues comerciales. & Community: Gratuito; Enterprise: Desde \textasciitilde\$1,000/mes. \\
\hline
\textbf{NanoMQ} \cite{zotero-item-232} & Broker ultraligero para edge computing, optimizado para IoT con MQTT 5.0. & Ligero y eficiente para edge computing, bajo consumo de recursos. & Menor ecosistema y soporte comunitario limitado. & Gratuito. \\
\hline
\textbf{AWS IoT Core} \cite{zotero-item-240} & Servicio en la nube de AWS para IoT, soporta MQTT 5.0, integración con AWS Lambda y reglas para datos. & Alta escalabilidad y seguridad, integración nativa con ecosistema AWS. & Dependencia de AWS y costos variables según uso. &Desde \$0.08/100,000 mensajes \$5-20/mes para proyectos pequeños. \\
\hline
\end{longtable}
\endgroup

